% !TEX root = ../Monografia.tex
% ---------------------------------------------------------------------------------------------------------------------------------------------------------
% CAPÍTULO - CONCLUSÕES
% ---------------------------------------------------------------------------------------------------------------------------------------------------------

\chapter{Conclusões}

% ---------------------------------------------------------------------------------------------------------------------------------------------------------
\section{Considerações Finais}
% ---------------------------------------------------------------------------------------------------------------------------------------------------------

Este trabalho partiu de uma observação empírica sobre o ecossistema Flutter: as ferramentas de diagnóstico de desempenho existentes são tecnicamente capazes, mas sua arquitetura externa ao processo de execução, dependente de configuração manual e desprovida de recomendações automáticas, impõe um atrito de uso incompatível com o ritmo iterativo do desenvolvimento moderno. A resposta proposta, o \texttt{collector\_flutter}, é um pacote Dart/Flutter que opera embutido na própria aplicação, coletando métricas de FPS, memória, rede, CPU, bateria e eventos customizados a cada ciclo de dois segundos, analisando-as por meio de heurísticas configuráveis e exibindo recomendações de otimização em um painel integrado, sem exigir configuração externa, conexão USB ou troca de contexto para uma ferramenta separada.

A arquitetura de três camadas (\textit{Data}, \textit{Domain} e \textit{Presentation}), baseada nos princípios da \textit{Clean Architecture}~\cite{martin2018clean}, demonstrou-se adequada para os objetivos do projeto. Na camada de dados, a coleta foi distribuída entre módulos de quadros, memória, rede, eventos, CPU e bateria. A lógica de análise permaneceu encapsulada em \texttt{Analyzer} e \texttt{Recommender}, enquanto a apresentação se apoiou principalmente em \texttt{DashboardPage} e \texttt{CollectorBloc}. Essa separação permitiu desenvolver, depurar e testar os componentes de forma independente. A extensibilidade da arquitetura foi confirmada na prática pela adição dos módulos de CPU e bateria sem modificações nas camadas superiores.

A validação funcional nos quatro cenários experimentais indicou a operacionalidade do sistema: foram observadas variações de \textit{frame time} durante \textit{rebuilds} intensivos, registro de requisições HTTP sob carga de rede, aumento de memória em leitura complementar por ADB e processamento simultâneo de múltiplos tipos de métricas sob carga combinada. As heurísticas de recomendação produziram sugestões compatíveis com os sintomas observados, especialmente nos cenários de rede e carga combinada.

Do ponto de vista acadêmico, o trabalho contribui para a discussão sobre instrumentação não-intrusiva de aplicações móveis, propondo uma implementação concreta em Dart com arquitetura documentada e avaliação funcional. Do ponto de vista prático, disponibiliza um pacote de código aberto que reduz a barreira de acesso ao diagnóstico de desempenho para desenvolvedores Flutter sem infraestrutura ou conhecimento especializado em \textit{profiling}.

% ---------------------------------------------------------------------------------------------------------------------------------------------------------
\section{Contribuições do Trabalho}
% ---------------------------------------------------------------------------------------------------------------------------------------------------------

As principais contribuições deste trabalho são:

\begin{enumerate}
  \item \textbf{Revisão crítica da literatura:} síntese de estudos sobre monitoramento de recursos em plataformas móveis, identificando a lacuna entre a disponibilidade de ferramentas de diagnóstico e sua adoção contínua durante o desenvolvimento, e fundamentando as escolhas de projeto do pacote.

  \item \textbf{Arquitetura modular documentada:} definição e implementação de uma arquitetura de três camadas para sistemas de monitoramento embutido em Flutter, com especificação formal de requisitos funcionais (RF1--RF5) e não funcionais (RNF1--RNF5) e mapeamento explícito entre requisitos, módulos e critérios de avaliação experimental.

  \item \textbf{Pacote funcional de código aberto:} implementação de doze módulos de coleta, análise e visualização disponibilizados no repositório público do projeto e também publicados no pub.dev, incluindo coleta de CPU via \textit{platform channel} Android e monitoramento de bateria, além de correção de cinco defeitos identificados durante o desenvolvimento.

  \item \textbf{Exportação de dados estruturada:} implementação do \texttt{ExportService} com exportação completa da sessão em JSON e de timings de frame em CSV, permitindo análise externa e reprodutibilidade das medições.

  \item \textbf{Suíte de testes automatizados:} conjunto composto por 60 casos de teste distribuídos em quatro arquivos, verificando cálculos internos, detecção de anomalias, geração de recomendações, serialização e comportamento de \textit{fallback} dos \textit{platform channels}. A lista de casos e o resultado da execução estão documentados no Apêndice~\ref{ap:testes}.

  \item \textbf{Validação funcional em quatro cenários controlados:} execução dos cenários de \textit{rebuilds} intensivos, carga de rede, alocação de memória e carga combinada, indicando operacionalidade do sistema e compatibilidade entre heurísticas e sintomas induzidos.

  \item \textbf{Discussão de implicações sociais e ambientais:} análise das dimensões de sustentabilidade digital, inclusão e responsabilidade ética associadas ao monitoramento de desempenho, situando a ferramenta no contexto mais amplo da engenharia de software sustentável.
\end{enumerate}

% ---------------------------------------------------------------------------------------------------------------------------------------------------------
\section{Limitações do Estudo}
% ---------------------------------------------------------------------------------------------------------------------------------------------------------

Apesar dos resultados obtidos, a avaliação realizada apresenta limites metodológicos que devem ser considerados na interpretação das conclusões:

\begin{itemize}
  \item \textbf{Natureza qualitativa da validação:} os cenários controlados permitiram verificar a coerência funcional das métricas e das recomendações, mas não constituem um experimento estatístico com amostra ampliada, repetições padronizadas e intervalos de confiança.

  \item \textbf{Foco principal em Android:} os testes funcionais foram conduzidos em ambiente Android, plataforma na qual os canais nativos de CPU e bateria foram efetivamente exercitados. No iOS, a leitura de bateria está disponível via \texttt{UIDevice.batteryLevel}, enquanto a leitura de CPU retorna valor sentinela em razão das restrições de \textit{sandbox}.

  \item \textbf{Inspeção informal da interface:} a clareza do painel foi analisada durante as execuções manuais pelo próprio pesquisador, observando legibilidade das métricas, distinção visual de severidade e associação entre alertas e cenários. Não houve aplicação de instrumento formal de coleta nem avaliação com desenvolvedores externos.

  \item \textbf{Comparação limitada com ferramentas de referência:} as leituras ADB serviram como apoio qualitativo. Em especial, \texttt{adb shell dumpsys gfxinfo} confirmou a superfície Flutter ativa, mas retornou contadores HWUI de frames zerados na execução em emulador, impossibilitando cálculo formal de erro de FPS em relação ao pacote.
\end{itemize}

% ---------------------------------------------------------------------------------------------------------------------------------------------------------
\section{Trabalhos Futuros}
% ---------------------------------------------------------------------------------------------------------------------------------------------------------

A continuidade deste trabalho pode avançar em cinco direções principais:

\begin{itemize}
  \item \textbf{Validação quantitativa formal:} repetir cada cenário em múltiplas rodadas controladas, com amostra suficiente para calcular médias, desvio padrão, intervalos de confiança e erro percentual em relação a ferramentas de referência.

  \item \textbf{Comparação sistemática com ferramentas especializadas:} executar os mesmos cenários com Android Profiler, Perfetto e comandos ADB, registrando séries temporais comparáveis de FPS, memória, CPU e eventos de renderização.

  \item \textbf{Ampliação de dispositivos:} avaliar o pacote em diferentes perfis de hardware, incluindo aparelhos de entrada, intermediários e topo de linha, além de versões variadas do Android.

  \item \textbf{Suporte iOS mais completo:} evoluir os \textit{platform channels} para métricas nativas de CPU e energia no ecossistema Apple, respeitando as restrições de \textit{sandbox} da plataforma.

  \item \textbf{Avaliação com desenvolvedores Flutter:} conduzir estudo de usabilidade com participantes externos para medir clareza do painel, utilidade das recomendações e facilidade de integração do pacote em projetos reais.
\end{itemize}

\clearpage
